%% main_ist_cut.tex --- Information and Software Technology (Elsevier) submission
%% Length-reduced version of main_ist.tex for the 15,000 word-equivalent limit.
%% Class: elsarticle, single-column preprint layout ("Your Paper Your Way").

\documentclass[preprint,12pt]{elsarticle}

\usepackage[T1]{fontenc}
\usepackage{booktabs}
\usepackage{graphicx}
\usepackage{amsmath}
\usepackage{amssymb}
\usepackage{xcolor}
\usepackage{url}
\usepackage{float}
\usepackage{placeins}
\usepackage{caption}
\usepackage{longtable}
\usepackage{tabularx}
\usepackage{array}
\usepackage{calc}
\usepackage{enumitem}
\usepackage{lineno}
\modulolinenumbers[5]

\providecommand{\tightlist}{\setlength{\itemsep}{0pt}\setlength{\parskip}{0pt}}
\setlist{itemsep=1pt, topsep=2pt, parsep=0pt, leftmargin=1.2em}
\setlength{\tabcolsep}{3pt}
\raggedbottom
\emergencystretch=3em
\hyphenation{re-pro-duc-tion spec-i-fi-ca-tions knowledge-graph-grounded
 taxonomy-grounded An-ten-na-Pod New-Pipe Thun-der-bird
 con-fi-dent cleanlab hi-er-ar-chi-cal cov-er-age IssueSpec}

\journal{Information and Software Technology}

\begin{document}

\begin{frontmatter}

\title{Typed Issue Specifications from App Reviews: What Scales, and a Downstream Benefit That Does Not Replicate}

\author[inst1]{Fabiha Jalal\corref{cor1}}
\ead{fabihajalal@iut-dhaka.edu}
\author[inst1]{Sadia Tasnim Dhruba}
\ead{sadiatasnim@iut-dhaka.edu}
\author[inst1]{Tazrian Zaman Tushti}
\ead{tazrian@iut-dhaka.edu}
\author[inst2]{Ajwad Abrar}
\ead{ajwadabrar@iut-dhaka.edu}
\author[inst2]{Md Kamrul Hasan}
\ead{hasank@iut-dhaka.edu}
\author[inst2]{Hasan Mahmud}
\ead{hasan@iut-dhaka.edu}

\cortext[cor1]{Corresponding author.}

\affiliation[inst1]{organization={Department of Computer Science, Islamic University of Technology (IUT)},
            city={Gazipur},
            country={Bangladesh}}
\affiliation[inst2]{organization={Systems and Software Lab (SSL), Department of Computer Science and Engineering, Islamic University of Technology (IUT)},
            city={Gazipur},
            country={Bangladesh}}

\begin{abstract}
\textbf{Context.} Developers receive millions of noisy app reviews, but defect trackers such as
JIRA and GitHub Issues cannot consume free text. They require structured fields such as issue
type, steps to reproduce, severity, and affected component. Review-mining research stops short of
that boundary, predicting labels or writing replies without producing the record a tracker can
ingest.

\textbf{Objective.} Can a typed intermediate representation (IR) of a review be produced at
corpus scale from noisy labels, and does supplying it to a response generator improve the reply a
user receives? We answer the first yes and the second no, and the second answer corrects a claim
we made ourselves.

\textbf{Method.} IssueSpec routes each review by issue class into one of five standards-body
templates, repairs the noisy labels the routing depends on with a verified-anchor procedure, and
feeds the result to a retrieval-augmented generator. We study 215,583 reviews from 58 Android
apps with a 5,230-review verified anchor, a 490-review three-rater gold standard, seven ablations,
and four language models. The downstream evaluation was run twice: first with conditions written
by different template composers, then with the generator held fixed so that only the IR varies,
all 200 blinded rows scored by three raters.

\textbf{Results.} Production works. The verified-anchor procedure lifts Cohen's $\kappa$ against
the human gold from 0.163 to 0.592, and holds at that level (0.616) on held-out reviews. Type
routing reaches 0.96 template-fill against 0.69 without it, a schema-conformance measure rather
than a quality one. The downstream benefit does not. The first round measured $+2.35$ Likert
points, but its arms differed in authorial register as much as in information, one emitting
lemmatised text at 78 words and the other polished prose at 123; the corrected comparison gives
$+0.04$ ($p = 0.54$), with raters disagreeing in direction. Three further components fail their
own tests: the knowledge-graph layer loses to count-matched flat clustering, our
extractive-coverage faithfulness scorer shows no rank correlation with human judgement, and the
constrained alignment layer never demonstrates the trade-off it was built to expose.

\textbf{Conclusions.} Typed issue specifications are producible from noisy reviews at corpus
scale, and that is the contribution we defend. Their claimed downstream payoff is not. We release
the arithmetic check that distinguishes our two rating rounds, because a structured intermediate
should be justified on what it produces rather than on an effect a fair comparison removes.
\end{abstract}

\begin{highlights}
\item A typed intermediate representation routes app reviews into five issue templates.
\item A verified-anchor procedure lifts Cohen's kappa from 0.163 to 0.592 at corpus scale.
\item LLM annotation noise concentrates on omitted classes, not on uniform per-item error.
\item A clean same-model re-run finds no downstream reply-quality gain from the typed spec.
\item Template-fill measures schema conformance, not the quality of the text in the fields.
\end{highlights}

\begin{keyword}
App review mining \sep Issue report generation \sep Intermediate representation \sep
Knowledge graph \sep Retrieval-augmented generation \sep Empirical software engineering
\end{keyword}

\end{frontmatter}

\linenumbers

\section{Introduction}\label{introduction}

App stores are where users talk back to developers, and the global stream runs into millions of
reviews a day~\cite{maalej2015bug, gao2019automating, dkabrowski2022analysing}. Those reviews
carry bug reports, feature requests, performance complaints, and regression notes as messy
free-form text, usually without the details required to reproduce a bug or scope a feature. Defect
trackers such as JIRA or GitHub Issues cannot consume raw text; they need structured fields such
as issue type, steps to reproduce, severity, and affected component~\cite{chaparro2017detecting,
zimmermann2010makes}, and filling those by hand does not scale. Years of work on automated review
pipelines has helped with classification, summarization, and response generation, but no system
produces the structured report a developer can paste into a tracker. D{\k a}browski et al.'s systematic review of 182 papers~\cite{dkabrowski2022analysing}
surveys this literature; it does not itself name this gap, and we state the gap as our own reading
of the work it covers. Section~\ref{related-work} treats each strand in turn: classifiers stop at
the label, response generators never write the defect out in structured form, agentic
software-engineering tools assume a well-formed issue as the agent's starting state,
knowledge-graph methods emit predictions and summaries a tracker cannot ingest, and constrained
RLHF targets general LLM safety rather than the domain-specific constraint family of developer
relations.

This work pursues three aims, stated as research questions.

\textbf{RQ1 (Translation Quality).} Can a typed IR, produced by a three-layer framework
(knowledge-graph construction, hierarchical clustering, standardized schema mapping), reach the
substantive completeness of human-written GitHub issues?

\textbf{RQ2 (Downstream Load-Bearing).} Does conditioning a response generator on the validated
IssueSpec measurably improve human-rated response quality over RAG alone?
Section~\ref{exp2-results} answers no, correcting our own earlier claim.

\textbf{RQ3 (Dual-Objective Alignment).} Can a CMDP-grounded RLHF loop through KTO, DPO, and
Constrained PPO jointly optimize response helpfulness and policy compliance for developer
relations? Section~\ref{exp3-results} answers no at the scale we can train, and reports what that
rules out.

We propose \textbf{IssueSpec}: a typed Intermediate Representation (IR) that converts each noisy
review into one of five standard report types, and a CMDP-grounded response generator that
consumes it under quality and compliance constraints (Figure~\ref{fig:architecture}).

\textbf{What is new, and what is assembled.} Several components are existing technique used as
intended, and we claim none of them: RoBERTa classification, confident
learning~\cite{northcutt2021confident}, Sentence-BERT with UMAP and HDBSCAN, PageRank,
retrieval-augmented generation, and Altman's CMDP formulation~\cite{altman2021constrained}. Three
things are new. \emph{First}, the typed intermediate representation: a review is routed by
predicted issue class into one of five standards-body templates, so the output is a typed record
with per-type obligations rather than a label or a summary, which no prior review-mining system
emits. \emph{Second}, the verified-anchor procedure, which makes that representation producible
from noisy LLM labels at corpus scale, and which works because the labeller omits two of its seven
classes entirely rather than because its errors concentrate on rare ones. \emph{Third}, a
correction and the check that produced it: we previously reported that the typed intermediate
carries downstream response quality, and Section~\ref{exp2-results} withdraws that, because the
comparison behind it varied the text generator as well as the specification, and because the
arithmetic of the two rating rounds turned out to distinguish independent human ratings from ones
derived by perturbing another rater's. We release both rounds and that check.

\textbf{The central claim, stated once.} The answer to which of five stages the paper is about is
Stage~3 and the object it produces, a \emph{typed, type-routed intermediate representation} of an
app review, together with the verified-anchor procedure that makes it producible at corpus scale
from noisy labels; everything else feeds that object or consumes it. Removing type routing drops
strict template-fill from 0.96 to 0.69 (Section~\ref{exp1-results}). We had also claimed a
downstream payoff, and Section~\ref{exp2-results} reports why we withdrew it: with the generator
held fixed, supplying the IR changes response quality by $+0.04$ Likert points ($p = 0.54$),
against the $+2.35$ a confounded design gave. Three further components fail their own tests, the
knowledge-graph layer against count-matched flat clustering, the \textsc{SpecCov} faithfulness
scorer against human judgement, and the CMDP alignment layer against the trade-off it was built to
expose. Alongside these we release a reusable check on rating provenance, described in
Section~\ref{exp2-results}. A reader who wants the claim and not the machinery can read
Sections~\ref{exp1-results}, \ref{exp2-results}, and~\ref{load-bearing} alone.
Table~\ref{tab:evidence} is the single map from component to claim to evidence to verdict, and
replaces the several overlapping summaries an earlier version carried.

\begin{table}[!htbp]\centering\scriptsize
\caption{Component-to-claim-to-evidence map, with the ablation that tests each component and the
verdict it returns. ``Removing this hurts'' is the only outcome that justifies a component;
several do not clear that bar and we report them as findings rather than as caveats. Ablation
identifiers (A1--A7) are defined in Section~\ref{ablations}; released artifacts are listed in
Section~\ref{reproducibility}.}\label{tab:evidence}
\begin{tabular}{@{}p{1.5cm}p{2.9cm}p{4.4cm}p{2.1cm}@{}}\toprule
component & what it must do & evidence, and the ablation that tests it & verdict \\\midrule
Stage 1 intake, verified anchor & repair a corpus label wrong 25\% of the time that omits two classes entirely, because Stage~3 routes on the type & $\kappa$ $0.163 \to 0.592$ against the three-rater gold, 0.616 held out (replay-dependent); macro F1 $0.22 \to 0.65$ with the two omitted classes recovered (\S\ref{exp1-results}) & \textbf{justified} \\
Stage 2 knowledge graph & group and rank 215{,}583 reviews so a developer has an order to work in & count-matched flat clustering beats the KG on all three intrinsic metrics (A1b) and on purity (A2); top-5 PageRank aspects cover 55\% of reviews & \textbf{prioritization yes, cluster quality no} \\
Stage 3 typed IR & emit typed tracker fields rather than prose & strict fill 0.96 vs 0.69 without routing, i.e.\ schema conformance and not quality (A3, Table~\ref{tab:stage3}); 5-dim rubric 3.94/5 on all 100 specs, LLM judge (Table~\ref{tab:rubric}) & \textbf{justified} \\
Stage 4 spec-aware RAG & supply house voice and precedent on top of the defect & retrieval moves automatic metrics $<1\%$ (A5); conditioning on the IR moves quality $+0.04$, $p=0.54$ (A6, \S\ref{exp2-results}), overturning the $+2.35$ we reported & \textbf{withdrawn / null} \\
Stage 5 CMDP RLHF & hold compliance as a constraint rather than a reward term & constraint never binds at distilGPT2 scale; with a policy able to violate it, binds on 13/90 steps and $\lambda$ rises to 0.65, with no measured tension ($\rho = -0.12$, $p = 0.25$), on a detector that scores hand-written violations as compliant (\S\ref{exp3-results}) & \textbf{partial testbed} \\
\textsc{SpecCov} scorer & rank specs the way a human judge of faithfulness does & no rank correlation with human faithfulness ($\rho=0.15$, $p=0.26$); orders generators backwards (\S\ref{speccov-validation}) & \textbf{withdrawn} \\
rating-provenance check & separate independent ratings from derived ones & the discarded round shows a single deviation magnitude and the reported round four (\S\ref{exp2-results}) & \textbf{released} \\
\bottomrule\end{tabular}\end{table}

 \begin{figure}[t]\centering
  \includegraphics[width=\textwidth]{fig_architecture.pdf}
  \caption{The IssueSpec five-stage pipeline. Reviews flow left to right through intake and
  classification (Stage~1), three-layer knowledge-graph clustering (Stage~2), LLM-agent IR
  translation (Stage~3), spec-aware RAG (Stage~4), and CMDP-grounded RLHF (Stage~5), each
  emitting a typed intermediate. Yellow checkpoints mark human-in-the-loop review under
  selective prediction~\cite{geifman2017selective, bansal2019beyond}; the red loop propagates Stage~5
  corrections upstream.}\label{fig:architecture}
  \end{figure}
%% =====================================================================
\section{Related Work}\label{related-work}

We organize related work by the five gaps of Section~\ref{introduction}, naming what each field
does, what assumption breaks for review-to-issue translation, and how IssueSpec differs.

\emph{App review classification and mining.} Maalej and Nabil~\cite{maalej2015bug} established a
four-class taxonomy (bug report, feature request, user experience, rating). The seven-class label
set we use is our own extension, adding \texttt{performance}, \texttt{usability}, and
\texttt{compatibility} and renaming the rating class \texttt{praise}; we are not aware of an
established seven-class app-review taxonomy, so the extension should be read as ours rather than
as canonical. Subsequent systems~\cite{chen2014ar, villarroel2016release, di2017surf,
kaur2023bert, shah2020mining, guzman2014users} refined accuracy and added aspect-level
granularity, and D{\k a}browski et al.~\cite{dkabrowski2022analysing} review 182 papers from 2012
to 2020, reporting median precision above 81\% and median recall around 83\%. An earlier version
of this paper attributed to that survey a claim it does not make, that label quality rather than
model capacity sets an F1 ceiling of 0.75 to 0.85: the survey reports precision and recall rather
than F1 and attributes limitations to evaluation-dataset size, so the label-quality argument is
ours and Section~\ref{exp1-results} is evidence for it rather than corroboration of theirs.
\emph{What breaks:} none of these systems emits a typed developer-routable artifact downstream of
the label.

\emph{Response generation and RAG.} RRGen~\cite{gao2019automating} and
CoRe~\cite{gao2021automating} pioneered end-to-end neural response generation on (review,
response) pairs without an explicit structured intermediate; vanilla RAG~\cite{lewis2020retrieval}
retrieves past responses in one pass and agentic variants~\cite{asai2024self,
khattab2022demonstrate} add critique-revise loops. \emph{What breaks:} none of these names the
underlying defect, and our $n=10$ feasibility study (Appendix~\ref{appendix:settings}) shows that
agentic iteration without a typed intermediate adds only a small gain on a rule-based scorer
(0.58$\to$0.70). We cannot claim the IssueSpec does better: once the generator is held fixed,
supplying it changes human-rated quality by $+0.04$ ($p = 0.54$, Section~\ref{exp2-results}), so
the contribution we defend is the specification itself rather than its effect on the reply.

\emph{Agentic software engineering.} SWE-Agent~\cite{yang2024swe},
RepairAgent~\cite{bouzenia2025repairagent}, and HyperAgent~\cite{phan2024hyperagent} consume
well-structured GitHub issues to produce code patches, and the SWE-bench
protocol~\cite{jimenez2023swe} makes the assumption explicit by pairing each instance with an
existing issue. \emph{What breaks:} the specification is the input contract for all of them, and
constructing one from noisy user-facing text is out of scope. IssueSpec is precisely the artifact
these systems are designed to consume.

\emph{Knowledge graphs for app reviews.} ReviewGraph~\cite{de2025reviewgraph} learns KG embeddings
for rating prediction, KGCPN~\cite{yu2024knowledge} produces extractive summaries, and Keertipati
et al.~\cite{keertipati2016approaches} apply centrality for coarse feature prioritization.
\emph{What breaks:} none maps KG nodes to a typed, developer-actionable schema. Our three-layer KG
pipeline is, to our knowledge, the first that emits a typed IR for review-to-issue translation.

\emph{RLHF and constrained optimization.} DPO~\cite{rafailov2023direct},
KTO~\cite{ethayarajh2024kto}, and PPO~\cite{schulman2017proximal} compress quality and compliance
into one scalar reward; Altman's CMDP framework~\cite{altman2021constrained} separates them, and
Safe-RLHF~\cite{dai2024safe} applies CMDP to general LLM safety with a domain-agnostic compliance
set. \emph{What breaks:} no prior work applies CMDP to dev-rel response generation, where the
compliance set is domain-specific (promise safety, information leakage, tone, legal safety). Our
Stage~5 formulation instantiates these four constraints inside the Lagrangian and releases a
proof-of-concept training loop.

Across these categories, existing systems occupy at most two of our five pipeline stages and none
produces a typed issue record. Composing all five with a standardized issue schema, three HITL
checkpoints, KG-centrality prioritization, and a CMDP-grounded RLHF formulation is the structural
contribution.

%% =====================================================================
\section{The IssueSpec Pipeline}\label{system-design}

The pipeline (Fig.~\ref{fig:architecture}) is a five-stage information-extraction
cascade~\cite{hearst1999untangling, sarawagi2008information}. Stages~1--3 produce the typed IR
from raw reviews, Stage~4 uses it to drive spec-aware response generation, and Stage~5 aligns the
generator via constrained RLHF. This section describes the framework only; evaluation design is
in Section~\ref{experiments}, exact settings in Appendix~\ref{appendix:settings}, results in
Section~\ref{results}, and the caveats each stage carries are stated with the results they
qualify. Table~\ref{tab:evidence} already gives the verdict on each stage, so a reader knows
before the description that two of the five do not justify themselves.

\subsection{Producing the IR}\label{aim1}

\textbf{Stage 1: Intake, Classification, and Aspect Extraction.} Stage~1 produces a multi-label
  class from a RoBERTa head fine-tuned on~\cite{maalej2015bug, guzman2014users}, aspect-level
  sentiment~\cite{guzman2014users} that keeps ``great UI but terrible battery'' as two pairs, and
  entity extraction for devices, OS versions, and named features, which Stage~3 needs as typed
  fields. The multi-label capacity is almost never exercised: on the 10{,}000-review generation
  subset only 5 reviews (0.05\%) receive more than one label, and the corpus-scale relabelling
  artifact stores a single argmax label per review, so the released runs are effectively
  single-label and readers should not credit the mixed-concern motivation with doing any work in
  our results. Low-confidence reviews route to HITL-1 by a fixed confidence threshold; letting a
  model abstain is a familiar idea~\cite{geifman2017selective}, but we implement a threshold
  rather than a calibrated selective-prediction procedure and report no risk-coverage behaviour,
  so the citation is background rather than a description of what runs.

  LLM annotation noise is the main challenge. On the 5{,}038-review praise stratum 25.8\% of LLM
  labels are wrong, and the labeller omits \texttt{compatibility} and \texttt{performance} almost
  entirely; per stratum the errors fall in the largest class, not the minority ones, so what a
  verified anchor repairs is the omission. We apply confident
  learning~\cite{northcutt2021confident} against an anchor of 5{,}230 expert-labeled reviews plus
  5{,}008 from Maalej~\cite{maalej2015bug}, gated by a confidence threshold on the anchor and a
  probability ceiling on the LLM label (Appendix~\ref{appendix:settings}). Compatibility needs a
  separate fix, since only 8 of 215{,}583 LLM labels carry it and Maalej has none: we seed it with
  200 template-generated synthetic reviews~\cite{joorabchi2013real} and 100 keyword-mined RRGen
  reviews. The lead author read both sets before training and corrected or
  dropped items that did not belong, but that pass was done by eye and no correction log was kept,
  so unlike the other annotation tasks in this paper it leaves no artifact a reader can check. We
  flag it because 300 of the 310 \texttt{compatibility} rows in V5's training set come from here,
  so the per-class F1 of 0.83 rests on data whose verification is asserted rather than evidenced.

  \textbf{Stage 2: Three-Layer Knowledge Graph.} Stage~2 builds, clusters, and prioritizes. The
  first layer links each review node to its aspects, with edges typed for sentiment and recency
  where available, but the released run populated less than that: the builder supports entity
  nodes and signed sentiment edges, yet the run behind our results passes empty entity lists and a
  constant neutral sentiment, leaving 8{,}404 review nodes, 10{,}534 aspect nodes, and no entity
  nodes. It is in practice the bipartite review-aspect graph PageRank runs over, and the sentiment
  and entity layers should be read as available in the implementation rather than exercised in the
  experiments. Aspects come from spaCy noun chunks with a curated keyword list and regular
  expressions, pruned by frequency and centrality; Hu and Liu~\cite{hu2004mining} is the ancestor
  of this approach rather than what we implement.

  The second layer clusters reviews first by aspect, then sub-clusters within each aspect using
  sentence embeddings~\cite{reimers2019sentence} and HDBSCAN~\cite{mcinnes2017hdbscan} directly on
  the 384-dimensional vectors. Earlier versions of this paper also named
  UMAP~\cite{mcinnes2018umap} here; that was wrong, since UMAP is used in the flat baseline this
  design is compared against and not in the released KG-hierarchical run, which applies no
  dimensionality reduction. Aspect-first keeps each sub-cluster on one theme, giving 605
  aspect-grounded sub-clusters of mean size 15.9 against 194 flat clusters of mean size 471.0. An
  earlier version of this paper reported 3.0 and 51.5; those were the number of representative
  reviews stored per cluster and a subsample artifact respectively, not cluster sizes. One scale
  caveat belongs here: the knowledge graph is built over a 10{,}000-review sample yielding
  8{,}404 usable reviews, so every KG-hierarchical number in this paper describes 3.9\% of the
  corpus, while the flat baseline it is set against clusters 91{,}368 reviews from the full
  corpus. The two Stage-2 designs never saw the same input, which is why we do not treat their
  intrinsic metrics as a head-to-head result (Section~\ref{cluster-validation-results}). The third
  layer prioritizes via PageRank~\cite{brin1998anatomy, page1999pagerank} on the bipartite
  review-aspect graph, and each cluster maps to the schema in Table~\ref{tab:templates} with a
  centrality-derived priority score.

  \textbf{Stage 3: Review-to-Issue Translation.} An LLM agent reads the standardized cluster
  schema and writes a specification using the type-routed templates in
  Table~\ref{tab:templates}: an LLM rather than rule-based extraction because reviews vary too
  much for fixed rules, and routed by issue type because a bug report needs reproduction steps, a
  feature request needs acceptance criteria, and a performance complaint needs latency or memory
  numbers, so a shared template would be incomplete or padded. The five templates are
  tracker-standard bug
  reports~\cite{zimmermann2010makes}, Cohn user stories~\cite{cohn2004user} with BDD acceptance
  criteria~\cite{hellesoy2017cucumber}, ISO/IEC~25010 non-functional
  categories~\cite{international2011systems}, Nielsen heuristics~\cite{nielsen1994usability}, and
  device-OS matrices~\cite{joorabchi2013real}. The output passes through HITL-2, where a domain
  expert scores each spec on a five-dimension rubric (completeness, accuracy, actionability,
  specificity, clarity, each 1--5), following the LLM-as-judge calibration convention
  of~\cite{gilardi2023chatgpt, pangakis2023automated, zheng2023judging}; five dimensions rather
  than one, because a single quality number conflates completeness with clarity and the LLM cannot
  tell which to fix on retry. Any spec failing on any dimension returns for a second pass with
  dimension-level feedback.

\textbf{\textsc{SpecCov} faithfulness metric.} We release \textsc{SpecCov}, an extractive-coverage
scorer over substantive-token overlap (length~$\geq 5$, stop-word filtered) with quoted-string and
coverage bonuses, following news-summarization conventions~\cite{grusky2018newsroom}; it targets
the extractive-coverage sub-type of faithfulness~\cite{maynez2020faithfulness}, a lexical
correlate of factual consistency~\cite{kryscinski2020evaluating} (other hallucination sub-types
in~\cite{ji2023survey} are out of scope). It is a lexical-grounding proxy, not a faithfulness
classifier: Section~\ref{speccov-validation} reports the experiment that withdraws it as one, and
that negative result only makes sense against this definition.

\begin{table}[t]
  \centering
  \footnotesize
  \caption{IssueSpec standardized issue schema. Each issue type follows the dominant industry or
  HCI standard. Each template is drawn from a published standard or reference and reproduces its
  named fields; we ran no conformance check against the standards documents themselves, so
  ``drawn from'' should not be read as ``certified against''.}\label{tab:templates}
  \begin{tabular}{@{}p{2.2cm}p{8.5cm}p{6.0cm}@{}}
  \toprule
  Issue type & Template (citation) & Justification \\
  \midrule
  \texttt{bug\_report} & title, severity, steps\_to\_reproduce, expected, actual, component &
  The tracker-standard field set in GitHub Issues and JIRA. Chaparro et
  al.~\cite{chaparro2017detecting} support three of the six (observed behaviour, expected
  behaviour, steps to reproduce). We previously attributed the whole set to
  Zimmermann et al.~\cite{zimmermann2010makes}; that was wrong twice over. They propose no
  template, and in their developer survey severity ranks last at 0\% importance and component at
  3\%, against 83\% for steps to reproduce. We keep both fields because trackers require them, not
  because developers report finding them useful \\
  \texttt{feature\_request} & Cohn user-story~\cite{cohn2004user} (As-a / I-want / So-that) +
  behavior-driven development (BDD) acceptance criteria~\cite{hellesoy2017cucumber} & Dominant
  industry format (Agile, Scrum, SAFe); BDD standardizes testable outcomes \\
  \texttt{performance} & Fields are \texttt{latency}, \texttt{startup\_time},
  \texttt{memory\_consumption}, \texttt{battery\_drain}, \texttt{network\_usage}, mapped onto
  ISO/IEC~25010~\cite{international2011systems} performance-efficiency sub-characteristics: the
  first two to time behaviour, the last three to resource utilization; capacity is not covered &
  International standard for software-quality NFR taxonomy. The field names are ours, chosen for
  mobile-app vocabulary, and are mapped to ISO categories rather than reproducing ISO terms \\
  \texttt{usability} & Nielsen-10 heuristics~\cite{nielsen1994usability} (visibility, match-real-world,
  user control, etc.) & Most-cited usability evaluation framework in human-computer interaction
  \\
  \texttt{compatibility} & Device-OS matrix (devices $\times$ OS
  versions)~\cite{joorabchi2013real} & Device/OS fragmentation is the canonical mobile-QA
  reporting axis \\
  \bottomrule
  \end{tabular}
  \end{table}

\subsection{Using the IR in Response Generation}\label{aim2}

Stage~4 is a RAG layer over five fixed sources, 15{,}100 documents in the released ChromaDB
index: changelogs and release notes (60 docs), FAQ and help documentation (40), 10{,}000 past
review-response pairs from RRGen~\cite{gao2019automating}, and 5{,}000 semantically similar past
responses, all pre-indexed, plus the validated IssueSpec injected per query. Together they give
the generator the answer, the voice, the closest precedent, and the structured intent.

\subsection{Aligning the IR-Aware Generator}\label{aim3}

Stage~5 frames alignment as a CMDP~\cite{altman2021constrained, dai2024safe},
  \[
  \max_\theta \mathbb{E}_{\pi_\theta}[R_{\text{quality}}] \quad
  \text{s.t.} \quad \mathbb{E}_{\pi_\theta}[C_{\text{compliance}}] \ge \tau,
  \]
  with the Lagrangian relaxation
  $\mathcal{L}(\theta,\lambda) = \mathbb{E}_{\pi_\theta}[R + \lambda(C - \tau)]$
  collapsing to single-objective optimization whenever the constraint is already satisfied. We use
  a CMDP rather than a scalar reward because a dev-rel response can be high-helpfulness and still
  out of compliance, like a polite promise the team cannot deliver, and a weighted sum hides the
  trade-off the constraint is meant to expose. The compliance scorer tracks four domain-specific
  dimensions absent from prior general-safety RLHF~\cite{dai2024safe}: promise safety, information
  leakage, tone, and legal safety. The trainer is Constrained PPO~\cite{schulman2017proximal} with
  a KL-regularized reference policy, with KTO~\cite{ethayarajh2024kto} and
  DPO~\cite{rafailov2023direct} as small-scale single-objective comparators
  (Appendix~\ref{appendix:settings}). HITL-3 supplies the two feedback streams, quality as reward
  and compliance as constraint, and also propagates backward, closing the Stage~5$\to$Stage~1 loop
  in Fig.~\ref{fig:architecture}.
%% =====================================================================
\section{Experimental Setup}\label{experiments}

\subsection{Dataset, Deduplication, and Annotation Provenance}
We use the RRGen corpus~\cite{gao2019automating}: 310{,}031 review-response pairs from 58 Android
applications. After exact-text deduplication that removes $94{,}448$ duplicate reviews
($-30.5\%$), 215{,}583 unique reviews form the working corpus, labeled across the seven classes.
Evaluation additionally uses a 490-review expert gold standard (70 per class $\times$ 7 classes)
and 100 stratified clusters for Stage~3.

Of the working corpus, 79.49\% of labels come from the V2 LLM annotator with no override, 18.08\%
are cleanlab-corrected against the verified anchor, and 2.43\% are directly human-verified. The
lead-author-versus-V2-LLM deviation is 25\% on the 5{,}230-row anchor (1{,}307 deviations,
24.99\%). We previously described this as replicated across two independent strata, the praise
stratum and the anchor, on a combined 10{,}271 rows. That was wrong: the praise stratum is 5{,}038
of the anchor's 5{,}230 rows, so it is 96\% of the same sample rather than a second one, and the
combined figure double-counted it. The rate is one measurement on one sample, and the praise
stratum alone deviates at 25.84\% (Table~\ref{tab:annot-deviation}). This $\approx 25\%$ deviation
rate, an order of magnitude higher than i.i.d.\ label-noise assumptions, motivates the
confident-learning procedure of Stage~1.

\subsection{Annotation Methodology}
Five kinds of annotator contribute, following the multi-annotator convention
of~\cite{gilardi2023chatgpt, pangakis2023automated}. The lead author verified 5{,}230 reviews
across minority strata, wrote 20 reference IssueSpecs, labeled the 490-review gold blinded to V2
LLM, and scored the Stage~4 sets with method IDs hidden until aggregation; two further human
raters annotated the same 490 items and, with her, formed the Stage~4 panel. A Qwen2.5-3B-Instruct
judge~\cite{zheng2023judging} scored IssueSpecs on the rubric and sub-clusters on the Y/P/N purity
rubric, calibrated against the lead author, and the corpus-labeling LLMs (V2, cleanlab-V2, V5) are
treated as independent raters for inter-rater statistics. The last two are model raters, not
people; three humans contribute in total.

\textbf{Label provenance, stated plainly.} No LLM-generated label is used as ground truth here.
LLM labels are the \emph{object under correction}, not the yardstick, and every reported accuracy
or $\kappa$ is measured against a human gold standard; where LLM raters do appear, in the
99-review classification check, they are labelled as such and no headline claim rests on them.

\textbf{How the gold was constructed, and what that costs.} Annotators did not label freely.
Each was shown the V5 classifier's predicted label and asked whether it was correct, supplying a
replacement when it was not, under a protocol telling them to accept a reasonable label and
overturn only when confident. That makes annotation fast and agreement high, and both flatter us.
It also couples the yardstick to the system being measured, since V5's label is the starting point
for every item and V2 and the cleanlab-corrected labels never received that benefit of the doubt.
The sample is stratified over V5's predicted classes as well, forcing 70 items per class into a
set where V2 predicts two of those classes almost never, so V2's per-class F1 of 0.000 on
\texttt{compatibility} and \texttt{performance} is partly a property of the sampling design. The
$\kappa$ trajectory should therefore be read as a comparison on a V5-seeded, class-balanced sample
rather than as a neutral benchmark, and a free-labelling gold would settle it. Conflicts are
resolved by majority vote: of 489 usable items 476 are unanimous and 13 are two-to-one splits,
with no three-way ties, and one item marked wrong by all three raters without a replacement is
excluded rather than guessed. All three first completed a shared 20-review calibration round
resolved against the written decision rules, and the per-item votes are released.

\textbf{A disclosure about an earlier annotation round.} The 490-review gold was annotated
twice. In the first round three sheets were distributed, and the two returned alongside the lead
author's carried annotation columns identical to hers cell for cell, so no genuine agreement could
be computed and they are not used anywhere in this paper; the round we report is the second, and
the agreement statistics come only from it. We disclose this for the same reason we disclose the
Stage-4 case in Section~\ref{exp2-results}: a paper offering a check for non-independent ratings
should say where its own artifact failed that check. Both rounds are released.

Stage~4 was likewise rated twice: a first 400-row round whose conditions came from template
composers, discarded for the reasons in Section~\ref{exp2-results}, and the 200-row round we
report, in which both arms come from the same model and differ only in whether the IssueSpec was
supplied, all rows scored independently by all three raters blind to condition. Ratings are
persisted as spreadsheets with JSON master keys, blinding shuffles conditions per pair, and the
assignment unseals only at aggregation. Table~\ref{tab:annot-deviation} consolidates agreement
across all pairs.

\begin{table}[t]\centering\footnotesize
  \caption{Annotator-vs-LLM and inter-annotator deviation. The $\approx 25\%$ deviation is
  measured on the human-verified anchor; the $\kappa$-progression reaches moderate
  agreement~\cite{landis1977measurement}; the V5 endorsement row is a consistency check, not independent validation, since V5 trained on those corrections;
  LLM-judge calibration $\Delta$ bounds judge bias. The $\alpha = 0.285$ row is an LLM-assisted
  panel, corrected from a mis-implemented 0.451 as described in the text.}\label{tab:annot-deviation}
  \begin{tabular}{@{}p{4.0cm}p{1.5cm}p{2.5cm}@{}}\toprule
  Annotator pair & sample & agreement / deviation \\\midrule
  Lead-author vs V2 LLM & 5{,}230 & 75\% / \textbf{25\%} mislabeled \\
  Lead vs V2 / cleanlab / V5 & 489 gold & $\kappa$: $0.163 \to 0.333 \to \mathbf{0.592}$ \\
  same, vs 3-rater majority gold & 489 gold & $\kappa$: $0.165 \to 0.334 \to \mathbf{0.590}$ \\
  3 human raters, Y/N verdict & 490 gold & Fleiss $\kappa = \mathbf{0.968}$; 479/490 unanimous on the verdict, 476/489 on the label \\
  3 human raters, 7-class label & 489 gold & Fleiss $\kappa = \mathbf{0.979}$; $\alpha = 0.979$ \\
  V2 LLM vs V5; cleanlab vs V5 & 215{,}583 & 71.96\% / 86.77\% agree \\
  cleanlab vs V5 (consistency, not independent) & 40{,}291 & 88.66\% support \\
  V5 vs gold, held out (replay-dependent, see \S\ref{exp1-results}) & 307 & $\kappa = \mathbf{0.616}$ \\
  3 human raters (Stage~4 quality, discarded round) & 400 rated rows & weighted $\kappa$ 0.970--0.986; within-1 100\% \\
  3 human raters (Stage~4 quality, reported round) & 200 rated rows & weighted $\kappa$ 0.03, 0.04, 0.50 \\
  3 human raters (Stage~4 helpful, discarded round) & 400 rated rows & $\kappa$ 0.678--0.836; $\alpha = 0.781$ \\
  3-rater \emph{LLM-assisted} classification check (lead + 2 Qwen prompts) & 99 reviews & Krippendorff $\alpha = 0.285$ \\
  Lead vs Qwen-judge (purity, rubric) & 50 / 28 & $\Delta = -0.035$ / $+0.23$ \\
  \bottomrule\end{tabular}\end{table}

Two three-rater statistics in Table~\ref{tab:annot-deviation} should not be conflated. The
classification gold uses three independent \emph{human} raters and reaches Fleiss $\kappa = 0.968$
on the Y/N verdict and $0.979$ on the seven-class label. The $\alpha = 0.285$ row is a separate
check on 99 reviews using the lead author plus two Qwen2.5-3B prompts as raters, measuring
seven-class classification agreement rather than cluster purity. We reported it as $\alpha =
0.451$ and as a purity check; both were wrong. The value came from an implementation that credited
each rating with a coincidence against itself, inflating observed agreement, and the corrected
figure agrees with the reference \texttt{krippendorff} package to four decimals; the three Cohen's
$\kappa$ in the same table were computed with scikit-learn and are unaffected. The correction
moves the row from the moderate band into the fair band, weakening the LLM-rater panel rather than
any claim we make.

\subsection{Experimental Design}
Because a single-LLM headline raises a structural-bias concern, Stage~3 is replicated across four
LLMs from three families, Claude Opus 4.7, Llama-3.3-70B-Instruct (via Groq), and
Qwen2.5-\{3B, 1.5B\}-Instruct (Table~\ref{tab:multi-llm}). We pre-register three experiments.

\textbf{Experiment~1 (RQ1, translation quality).} Five conditions on the same 100 stratified
clusters: (a)~LLM with taxonomy grounding (full Stage~3), (b)~free-form LLM, (c)~raw concatenation
as a lower bound, (d)~lead-author reference, and (e)~GitHub issues mined from three open-source
Android projects as the practitioner upper bound, each scored on the rubric with paired Wilcoxon
tests.

\textbf{Experiment~2 (RQ2), as reported.} The design we report holds the generator fixed. The
same model, given the same retrieved context, writes two replies for each of 100 reviews,
differing only in whether the validated IssueSpec is in the prompt. All 200 rows are rated blind
by three raters on quality, specificity, and helpfulness, and analysed with a paired Wilcoxon test
and a paired bootstrap interval. \emph{The discarded design}, described because the paper reports
its withdrawal rather than because it supports a claim, used the same 100 clusters to drive four
conditions produced by template composers rather than by a model: (a)~an rrgen-style baseline
(review only), (b)~a prompt baseline (CoRe-style dev-rel prompt with lightweight entity
extraction), (c)~RAG without IssueSpec, and (d)~the full system. It was evaluated with automatic
metrics (BLEU, ROUGE-L, BERTScore) and 400 blinded paired human ratings, with significance by
paired Wilcoxon, Friedman-Nemenyi~\cite{demvsar2006statistical}, Bradley-Terry, and McNemar at
Bonferroni-corrected $\alpha = 0.05$, and effect sizes following~\cite{romano2006appropriate,
cliff1993dominance}. Both rounds run single-shot vanilla RAG~\cite{lewis2020retrieval} rather than
an agentic loop, so the IssueSpec's contribution is not confounded with orchestration; the agentic
comparison against~\cite{asai2024self, khattab2022demonstrate} runs on a 10-review feasibility
sample and is a demonstration rather than a powered comparison.

\textbf{Experiment~3 (RQ3, dual- vs.\ single-objective RLHF).} Five policies run on a
  distilGPT2 proof-of-concept (82M parameters, 400 SFT samples, 296 KTO pairs, 100 DPO pairs):
  SFT-base, single-objective KTO~\cite{ethayarajh2024kto}, single-objective
  DPO~\cite{rafailov2023direct}, a constrained-proxy variant, and dual-objective Lagrangian
  Constrained PPO~\cite{schulman2017proximal, dai2024safe}. We measure head-to-head BLEU, ROUGE-L,
  and BERTScore on a 100-review test set, plus the rule-based safety-violation rate.
%% =====================================================================
\section{Results}\label{results}

Each research question owns one subsection: Section~\ref{exp1-results} answers RQ1,
Section~\ref{exp2-results} answers RQ2 and explains what our earlier positive answer was
measuring, and Section~\ref{exp3-results} answers RQ3 in the negative.
Sections~\ref{cluster-validation-results} and~\ref{ablations} test individual components, and
Section~\ref{load-bearing} states which earn their place.

\subsection{Is the IR Producible at Scale?}\label{exp1-results}

Under strict content-validity criteria (three or more action-verb reproduction steps,
descriptions $\geq 30$ words, a full As-a/I-want/So-that triple), the LLM-with-taxonomy condition
reaches a template-fill rate of 0.96 (Table~\ref{tab:stage3}), free-form LLM 0.69, the same as the
lead-author reference set ($n=20$), raw concatenation 0.34, and mined GitHub issues 0.53. The
LLM-vs-GitHub gap reflects coverage rather than content quality, since the LLM is prompted to fill
every field whereas humans leave optional fields blank, and a very-strict sensitivity sweep
preserves the ordering, Claude dropping from 0.96 to 0.71 with that gap shifting only $-0.085$.

The limit of this measure is worth stating exactly. The free-form condition fills the four shared
fields and none of the type-specific ones, so its score is fixed by the issue-type mix, $(30 \cdot
4/7 + 30 \cdot 4/6 + 40 \cdot 4/5)/100 = 0.691$, with zero variance across specs: the 0.69 is a
constant, not a measurement, and the gap says only that routing produces type-specific fields at
all. Scoring the same conditions on bare field presence gives 1.000 against 0.678, versus 0.959
and 0.678 under the strict criteria, so for the free-form arm the content thresholds remove
nothing and the whole reported gap is schema conformance. A hand-written bug specification whose
description is ``something'' repeated thirty-five times, with three five-word steps built from the
criterion's own action-verb list, passes all seven fields and scores 1.000, above our headline. A
reader should treat 0.96 as schema conformance and nothing more.

\textbf{\textsc{SpecCov} scores, and a correction to how they were computed.} An earlier
version of our scorer applied per-condition floors, raising \texttt{raw\_summary} to 5 and the
human reference to 4 regardless of content. That is not a measurement, because a score whose value
depends on the label of the thing being scored cannot compare labels, so we removed the floors.
Across 320 specs the LLM-with-taxonomy condition averages 4.19, free-form LLM 3.38, raw
concatenation 4.47, and the human reference 3.45. The two LLM conditions are unchanged, because no
floor was ever applied to them; the floors had inflated the two comparison conditions, so removing
them makes our own system look better rather than worse, and we say so to put the direction of the
correction on the record. On their own these numbers rank raw concatenation highest, the trivial
upper bound obtained by copying the source, and Section~\ref{speccov-validation} shows they should
not be read on their own. Top-ranked aspects by PageRank centrality are \texttt{ad} (0.040),
\texttt{phone} (0.011), and \texttt{ads} (0.010), with \texttt{edit} (0.009) and \texttt{battery}
(0.007) next; an earlier version listed \texttt{battery} third, skipping the two above it.

\subsubsection*{SpecCov does not measure faithfulness}\label{speccov-validation}
An extractive-coverage score is only useful if it agrees with what a person calls faithful. Three
human raters independently scored 60 specs, 20 from each of three generators, shuffled and
unlabelled, on a 1--5 faithfulness scale, and separately listed every claim the source reviews do
not support. One qualification first: one rater completed the corrected sheet, which shows exactly
the evidence \textsc{SpecCov} scores against, while the other two completed an earlier draft
showing a subset of it, so their flags are biased toward over-flagging and we lead with the rater
who saw the full evidence. The metric fails all three checks on either reading. \emph{It does not
rank specs the way humans do:} Spearman $\rho$ between \textsc{SpecCov} and human faithfulness is
$+0.045$ ($p = 0.73$), $+0.015$ ($p = 0.91$), and $+0.226$ ($p = 0.08$) for the three raters, and
$+0.147$ ($p = 0.26$) against their mean, none significant. \emph{It does not detect the claims
humans call unsupported:} one rater marked 45 of the 60 specs as containing at least one
unsupported detail and 15 as clean, and \textsc{SpecCov} scores those groups at 3.64 and 3.93, a
difference in the right direction but far from significant (Mann-Whitney $p = 0.34$), where a
metric that detected hallucination would separate them. \emph{It orders the generators
backwards:} it ranks the taxonomy condition (4.20 on this subsample) above the human reference
(3.45) where the humans rank them the other way, reference 4.10 and taxonomy 3.50, because
templated specs reuse source vocabulary densely, which the score rewards, while a human writer
condenses and paraphrases, which it penalizes even when nothing is invented.

We therefore withdraw \textsc{SpecCov} as a faithfulness measure. Lexical grounding is not a
usable proxy for faithfulness on this task, and the automatic 320-spec rubric that appeared to
support it is not independent evidence, since its faithfulness dimension is computed by the same
overlap procedure. Entailment-based measurement~\cite{laban2022summac, honovich2022true,
durmus2020feqa} is the direction we would now take. We release the scorer, the 60-spec human
ratings, and this analysis, because a documented negative result on a plausible-looking metric
is more useful than a metric that was never tested.

\textbf{Cross-LLM replication (bias control).} Llama-3.3-70B reaches 0.80 template-fill,
Qwen2.5-3B 0.74, and Qwen2.5-1.5B 0.48 under the same criteria, and the
templated~$>$~free-form~$>$~raw ordering holds across all four (Table~\ref{tab:multi-llm}, with
the two caveats its caption records, one a withdrawn row). On two per-class structural checks Llama-3.3-70B attains complete template compliance on the
stratified sample, all 29 bug specs carrying three or more action-verb reproduction steps (vs.\
Claude's 73\%) and all 30 feature specs the As-a/I-want/So-that triple (vs.\ 97\%); these measure
structural compliance, not content correctness. What the bias control shows is cross-family rank
preservation: the IssueSpec contract transports across families rather than depending on one
proprietary generator.

\textbf{Cross-family rubric quality.} Template-fill says nothing about content quality, so we ran
the same Qwen2.5-3B judge and rubric over the other generators' specs on matched clusters. On the
14 clusters all four cover (bug reports only, where the two Qwen runs overlap), mean rubric score
follows the template-fill order: Claude 4.09, Llama-3.3-70B 3.96, Qwen2.5-3B 3.83, Qwen2.5-1.5B
3.77. At $n=14$, or $n=13$ for the three pairs involving Qwen2.5-1.5B, none of the six pairwise
differences reaches significance (Wilcoxon $p$ 0.18 to 0.75), and the adjacent gaps are the size
of the run-to-run generation noise of Section~\ref{reproducibility}, so the ordering is consistent
with the template-fill result rather than independent confirmation of it. A wider
Claude-versus-Llama comparison on 28 clusters spanning all five issue types gives 4.01 against
3.64, a nominal $+0.38$ at $p = 0.015$, which we previously asserted as the one cross-family
quality difference we stand behind. We withdraw that assertion: it belongs to a family of seven
Wilcoxon tests with the same judge and rubric, Bonferroni or Holm correction puts it at $p =
0.104$, and an exact sign test on the same 28 pairs gives $p = 0.078$. We no longer claim a
cross-family quality difference at all. The rubric spread is much narrower than the (unmatched)
template-fill spread, which is itself informative: the capability gap shows up mostly as whether
fields get filled, with completeness separating the models (4.00 down to 3.00) and accuracy nearly
flat (4.50 for three of four). One further caveat: the judge is itself Qwen2.5-3B-Instruct, so
that row carries a self-preference risk the others do not.

\begin{table}[t]\centering\footnotesize
\caption{Multi-LLM evidence summary across three model families (Anthropic, Meta, Alibaba).
Cross-family rank preservation in Stage~3 is the structural-bias control; LLM-as-judge
calibration $\Delta$ vs the lead author bounds judge bias. \emph{Two rows carry caveats a reader
needs.} The Stage-4 row previously reported 0.61 versus 0.54 for ``Claude'' against Qwen-3B. We
withdraw it: the file labelled Claude is \texttt{responses\_rrgen\_baseline.json}, which our own
release notes record as the output of a deterministic template composer rather than of any
language model, so the row compared a template against a generation. This is the same confound we
withdrew in Experiment~2 and we had not noticed it here. The Stage-3 row is a real cross-family
comparison but is not matched: Claude is scored on 100 clusters and Llama on 99, both spanning
five issue types, whereas both Qwen conditions are 15 clusters of \texttt{bug\_report} only. The
0.96-versus-0.48 span therefore compares a five-type mean against a single-type mean and should be
read as an ordering, not as a magnitude.}\label{tab:multi-llm}
\begin{tabular}{@{}p{2.7cm}p{2.9cm}p{2.4cm}@{}}\toprule
Use & Model (family) & Result \\\midrule
Stage~3 generate & Claude (Anthropic) / Llama-3.3-70B (Meta) / Qwen-3B (Alibaba) / Qwen-1.5B (Alibaba) & 0.96 / 0.80 / 0.74 / 0.48 \\
Stage~4 generate & \emph{withdrawn} (see note) & \emph{withdrawn} \\
5-dim rubric judge & Qwen-3B, all 100 specs & 3.94/5 (28-spec subsample: 3.89) \\
Cross-family rubric & Claude / Llama / Qwen-3B / Qwen-1.5B & 4.09 / 3.96 / 3.83 / 3.77 (\S\ref{exp1-results}) \\
Y/P/N purity judge & Qwen-3B vs lead-author & 0.625 vs 0.66; $\Delta=-0.035$ \\
\bottomrule\end{tabular}\end{table}

\textbf{Five-Dimension Rubric.} The rubric~\cite{gilardi2023chatgpt, pangakis2023automated,
zheng2023judging} is operationalized with Qwen2.5-3B-Instruct as judge over the complete set: all
100 taxonomy specs and all 20 lead-author reference specs in one run (Table~\ref{tab:rubric}). The
taxonomy condition reaches $3.94/5$ (sd 0.59), above the hypothesized lower bound of 3.5, and an
earlier stratified 28-spec subsample gave $3.89$, so subsampling cost $0.05$ of a Likert point.
Per issue type it ranges from 4.23 on performance down to 3.45 on usability. The reference set is
a judge calibration, not a competitor: it scores $3.45$, with the gap concentrated in
actionability (2.85 versus 3.86), because the lead author left template-specific fields empty,
writing a reference spec by hand being a different task from filling a routed template. The
comparison therefore says that a human writing freely produces less tracker-ready output than a
type-routed generator. It is not evidence that the specs are better than expert judgement on
content, and we do not read it that way.

\begin{table}[t]\centering\footnotesize
\caption{Experiment~1 five-dimension rubric (Qwen-judge, 1--5 scale) on the \emph{complete}
spec set: all 100 taxonomy specs and all 20 lead-author reference specs, scored in one run so
the two columns are directly comparable. The middle column repeats the earlier stratified
28-spec subsample under the same judge and rubric; the gap between it and the full set is
$+0.05$, which is the subsample error. The LLM-as-judge protocol follows the
calibration-against-reference convention of~\cite{gilardi2023chatgpt, pangakis2023automated,
zheng2023judging}.}\label{tab:rubric}
\begin{tabular}{lrrr}\toprule
dimension & taxonomy, $n=100$ & taxonomy, $n=28$ & lead-author ref, $n=20$ \\\midrule
completeness & \textbf{3.72} & 3.79 & 3.20 \\
accuracy & \textbf{4.43} & 4.39 & 4.10 \\
actionability & \textbf{3.86} & 3.71 & 2.85 \\
specificity & \textbf{4.15} & 4.07 & 4.05 \\
clarity & \textbf{3.52} & 3.50 & 3.05 \\\midrule
overall mean & \textbf{3.94} & 3.89 & 3.45 \\
sd across specs & 0.59 & -- & 0.41 \\
\bottomrule\end{tabular}\end{table}

\begin{table}[t]\centering\footnotesize
  \caption{Experiment~1 substantive template-fill (strict criteria). The 0.96 vs 0.53
  LLM-vs-GitHub gap is coverage, not content quality; the strict-vs-very-strict sweep preserves
  ordering. The exact-zero cells in \texttt{bugs:repro} and \texttt{features:story} columns for
  (b), (c), (d), and (e) are structural: free-form LLM, raw concatenation, lead-author, and
  human-GitHub specs do not emit the template-specific reproduction-step list or
  As-a/I-want/So-that triple at all when not prompted for them, so the strict-fill rate is zero
  by construction.}\label{tab:stage3}
  \begin{tabular}{lrrr}
  \toprule
  condition & strict & bugs:repro & features:story \\
  \midrule
  (a) LLM + taxonomy & \textbf{0.96} & \textbf{0.73} & \textbf{0.97} \\
  (b) LLM free-form & 0.69 & 0.00 & 0.00 \\
  (c) raw concat & 0.34 & 0.00 & 0.00 \\
  (d) lead-author ref ($n=20$) & 0.69 & 0.00 & 0.00 \\
  (e) human GitHub (3 repos) & 0.53 & 0.13 & 0.00 \\
  \bottomrule
  \end{tabular}
  \end{table}

\textbf{Stage~1 Classification Recovery.} The Cohen's $\kappa$ trajectory against the 490-review
expert gold crosses three Landis-Koch~\cite{landis1977measurement} agreement bands
(Table~\ref{tab:annot-deviation}): the V2 LLM achieves $\kappa = 0.163$ (slight); after cleanlab
correction~\cite{northcutt2021confident} $\kappa = 0.333$ (fair); the V5 classifier retrained on
the corrected labels reaches $\kappa = 0.592$ (moderate). The trajectory does not depend on any
single annotator: scored against the majority-vote gold of the three-rater human panel ($n = 489$
after excluding one review left without a correction label), the same three classifiers reach
$\kappa = 0.165$, $0.334$, and $0.590$, within $0.002$ at every step. On the held-out 307-review
subset (183 of the 490 entered training, mostly through the stratified cap rather than the anchor
proper), V5 reaches $\kappa = 0.616$. We previously wrote that this is ``slightly above the 0.592
on the full 490, ruling out a contamination artifact.'' Two audits make that untenable as written.

First, ``held out'' is not exactly true. All 100 keyword-mined \texttt{compatibility} rows in the
training augmentation are verbatim copies of reviews already in the corpus, appended past index
215{,}583 with the label changed. Eight of them duplicate a held-out review, seven with the same
label the expert assigned, and the index-based replay cannot see this because it filters on
position rather than text. Dropping those eight gives $\kappa = 0.609$ on $n = 298$.

Second, and more seriously, the split is a replay rather than a recorded fact, and 0.616 is the
most favourable replay. The reconstruction shuffles each class with seed 42 and iterates classes
in Python dictionary insertion order. Iterating them in sorted order instead, changing nothing
else, gives $\kappa = 0.594$; seeds 0, 1, 2, 7, 123 and 2024 give 0.602, 0.605, 0.617, 0.590,
0.583 and 0.606. The published figure is at the top of that range, and the direction of the
inequality against 0.592 flips under most alternatives. What survives is the weaker statement:
held-out $\kappa$ is indistinguishable from full-set $\kappa$, across replays and after removing
the duplicated rows. That still argues against memorisation, which is what the sentence was for,
but it does not do so by being larger, and we should not have leaned on a direction this fragile.

The per-class recovery makes the structural collapse visible. Macro F1 against the gold rises
from 0.22 (V2) through 0.38 (cleanlab-V2) to 0.65 (V5), and on the two minority classes the V2 LLM
cannot see at all, F1 goes from 0.00 to 0.83 for \texttt{compatibility} and 0.00 to 0.77 for
\texttt{performance}, the compatibility figure resting on the 300 augmented rows discussed below. Applied across the full corpus, V5
agrees with the cleanlab correction on 88.66\% of the 40{,}291 corrected cases. V5 was trained on
a corpus in which 38{,}985 of those corrected rows carry the corrected label, so its agreement is
substantially a measure of fit to its own training signal, not an independent second opinion. We
report it as a consistency check and do not claim the two signals are independent. The choice of
confident-learning thresholds is not load-bearing: a $3\times 3$ grid over $(\text{anchor\_conf},
\text{max\_llm\_prob})$ keeps V5 endorsement in an 82.7--91.0\% band, with the headline $(0.70,
0.20)$ at 88.66\% near the grid center.

The held-out split itself needs describing, because it is not unbiased in the way we first said.
V5's training corpus uses a stratified per-class cap of 15{,}000 rows at seed 42, and replaying
that cap leaves 307 of the 490 gold reviews unseen, 306 of them scorable. The cap binds only on
the four classes V5 predicts most often, so membership of the unseen set is decided by V5's
\emph{predicted} class rather than at random. Scored by expert label the unseen 307 still spans
all seven classes, including 38 \texttt{compatibility}, 32 \texttt{usability}, and 31
\texttt{performance} reviews, so it does cover the classes carrying the recovery; what it cannot
rule out is that reviews V5 predicted confidently into large classes, which may be easier, are
over-represented there. We release the replay script, since the split is a faithful
reconstruction rather than a recorded fact.

To rule out memorization of the 200 synthetic compatibility templates used during V5 training, we
evaluate V5 against Maalej's 5{,}008 labels under a label-space cross-protocol test: predictions
are scored against Maalej's taxonomy, which has no \texttt{compatibility} class, so any
compatibility prediction is out-of-distribution with respect to the evaluation labels. On the five
overlapping classes ($n=4{,}560$), V5 achieves macro F1~=~0.676, weighted F1~=~0.730, and
accuracy~=~72.0\%. V5 also flags 432 reviews (8.6\%) as \texttt{compatibility}; the
top-twenty-confidence flags read as canonical compatibility complaints (``crashes on iphone 5 and
ipad 3''), all labeled \texttt{bug\_report} in Maalej because Maalej offers no alternative, so on
this evaluation the detector performs concept-recognition rather than template-matching.

A caveat about the class itself, found by auditing the training file. Of the 310
\texttt{compatibility} rows V5 trains on, 300 are augmentation. The corpus is lowercased and
lemmatised throughout, and not one of its 215{,}583 rows contains a single uppercase character,
whereas 194 of the 200 synthetic rows do. V5 is a cased RoBERTa, so any capital letter is a free
and perfect marker for ``synthetic compatibility''. The internal test F1 of 0.74 for this class is
therefore not a meaningful number and we withdraw it. The gold-standard compatibility F1 of 0.83
is not affected by the cue, because the gold is lowercase corpus text where it never fires; it is
carried instead by the 100 mined rows, which are the verbatim duplicates described above. The
cross-protocol result stands, since Maalej's text is also not from the synthetic set, but a reader
should treat every compatibility number in this paper as resting on 300 augmented rows rather than
on natural corpus evidence.

\subsection{Is the IR Load-Bearing? A Result We Had to Withdraw}\label{exp2-results}

An earlier version of this experiment reported that adding the IssueSpec to the response
generator raised human-rated quality by $2.35$ Likert points. That number was wrong, not in its
arithmetic but in what it measured, and the correction is the most important thing this section
has to report.

\textbf{What was wrong.} The four conditions in that experiment were not produced by a language
model but by deterministic template composers, Python scripts selecting sentence fragments from a
hand-written bank, and those composers differed in authorial register as much as in what
information they were given. The \texttt{no\_spec} arm emitted lemmatised, ungrammatical text
averaging 78 words (``we be sorry that the login loop you describe have got in the way''); the
\texttt{full} arm emitted polished prose averaging 123 words, with cleanup passes and an empathy
phrasebook no other arm received. A rater comparing them was separating two writing styles, and
the presence of the IssueSpec was confounded with that separation.

\textbf{The clean comparison, and what we can and cannot verify about it.} We re-ran the
experiment with the generator held fixed: both arms come from the same model on the same 100
reviews, differing in whether the validated IssueSpec appears in the prompt. One premise of that
design is not evidenced by our release and we will not assert it. Neither response file stores the
prompt or the retrieved text, and the two record their retrieval sources under different label
schemes with different counts, so no released row has an identical recorded source list; both arms
draw from the same candidate pool, but whether the retrieved context was identical per review is
not recoverable. A reader should treat this as holding the generator and review set fixed rather
than as a strict single-variable manipulation. Retrieval differences would add noise rather than a
systematic advantage to either arm, which makes a null easier to obtain, so this does not rescue
the withdrawn positive result; it does mean the null is weaker evidence than a fully controlled
design would give. The two arms average 105 and 97 words with distinct-1 at 0.736 and 0.760, close
on the surface features a rater could use to tell them apart, though the length difference is
itself significant (paired Wilcoxon $p = 6 \times 10^{-8}$).

\textbf{The IssueSpec does not measurably improve response quality.} Pooled over the three raters
the gain is $+0.04$ Likert points, 95\% bootstrap CI $[-0.10, +0.18]$, paired Wilcoxon $p = 0.54$
on 100 pairs, Cliff's $\delta = 0.04$. The two replies per review are shown in a randomised order;
the slot assignment came out 64/36 rather than balanced and raters show a small first-slot
preference, so we also computed a slot-stratified estimate, which moves the pooled figure to
$+0.03$ and each rater's by at most 0.03. We use the unadjusted numbers throughout so that
Table~\ref{tab:stage4-rerun} is on one estimator, and release the stratified values with the data.
Specificity behaves the same way, pooled $-0.00$ ($p = 0.87$), and the helpful rate is flat at
94\% in both arms. On this design we cannot claim that supplying the typed specification improves
the response a user receives, and we should be equally careful in the other direction: with 100
pairs and the observed spread the smallest detectable effect at 80\% power is about a fifth of a
Likert point, so we have ruled out a large effect and not a small one, and detecting a gain the
size of the one we measured would need on the order of 2{,}500 pairs. The effect, if any, is small
enough that this study cannot see it, not zero.

\textbf{The raters disagree with each other, and that is the interesting part.}
Table~\ref{tab:stage4-rerun} breaks the result out. The lead author, who designed the system,
prefers the IssueSpec arm decisively: $+0.56$ quality ($p = 2 \times 10^{-11}$, 59 wins, 5 losses,
36 ties) and $+0.39$ specificity ($p < 10^{-7}$). The two independent raters lean the other way,
$-0.29$ and $-0.15$, neither significant at $\alpha = 0.05$ under scipy's default zero-handling
($p = 0.06$ and $0.09$; the first has 45 tied rows of 100, and under Pratt's method that becomes
$p = 0.044$, while its bootstrap interval $[-0.56, -0.01]$ excludes zero). We state both rather
than pick the one that reads better: on that rater the evidence sits exactly on $\alpha = 0.05$
and moves either side depending on defensible choices. The disagreement between raters is larger
than the effect under measurement, and inter-rater reliability is correspondingly low,
quadratic-weighted $\kappa$ of 0.03, 0.04 and 0.50 against 0.97 in the discarded template-based
round, where high agreement was a symptom of the confound: when one arm is visibly broken, raters
agree easily.

We looked for a content mechanism behind the lead author's preference and did not find one.
Correlating each rater's per-review preference against the spec-derived features the IssueSpec
should introduce, named devices, reproduction steps, and component mentions, produces no
correlation that survives inspection; the largest is $\rho = 0.19$ ($p = 0.06$) between the lead
author's preference and reply length. Two readings remain open, that she perceives triage value a
general reader has no reason to weigh, in which case the right evaluation asks developers, or that
the preference is expectancy on the part of the person who built the system. This design cannot
separate them, and we assert neither.

\begin{table}[t]\centering\footnotesize
\caption{Stage-4 re-run, same model with and without the IssueSpec, 100 paired reviews, three
raters, all 200 rows scored by each. All figures are unadjusted, so the gain column equals the difference of the two condition
columns; slot-stratified values are in the released JSON and differ by at most 0.03. The pooled
effect is null, and the disagreement between the lead author and the two independent raters is
larger than the effect being measured.}\label{tab:stage4-rerun}
\begin{tabular}{@{}lrrrrr@{}}\toprule
rater & no\_spec & full & gain & 95\% CI & Wilcoxon $p$ \\\midrule
lead author & 4.04 & 4.60 & $+0.56$ & $[+0.44, +0.68]$ & $2 \times 10^{-11}$ \\
independent rater 1 & 3.87 & 3.58 & $-0.29$ & $[-0.56, -0.01]$ & 0.06 \\
independent rater 2 & 3.57 & 3.42 & $-0.15$ & $[-0.33, +0.03]$ & 0.09 \\\midrule
\textbf{pooled} & 3.83 & 3.87 & $\mathbf{+0.04}$ & $[-0.10, +0.18]$ & \textbf{0.54} \\
\bottomrule\end{tabular}\end{table}

\textbf{Why we believe the raters this time, and what the check does and does not show.} In the
discarded round the two additional raters deviated from the lead author on only 16 and 19 of 400
rows, those deviation sets did not overlap on a single row, and every deviation was exactly one
point. We initially read the zero overlap as the damning fact. It is not: with deviations that
rare, the hypergeometric probability of zero overlap under independence is 0.45, so it is the
single most likely outcome and carries almost no information. The discriminating observation is
the second: on the quality column every deviation in both sheets had the same magnitude, one
point, which is what a vector perturbed by a fixed step produces and not what human disagreement
produces. We name the column because the specificity column in the same sheets does show
magnitudes of one, two and three, so the signature is present in one column and not the other. In
the reported round the raters differ from the lead author on 133 and 145 of 200 rows, the
deviation sets overlap on 101 where chance predicts about 96, and deviations take four distinct
magnitudes. We release the check with the data and state its weak prong plainly, because a check
offered as a contribution should come with its own limits attached.

\textbf{What survives.} Less than we first thought. The claim that the typed IR carries
downstream response quality does not survive, and neither does the retrieval finding we treated as
the surviving half. We first reported that retrieval alone was indistinguishable from a no-context
baseline, then corrected that to significantly worse; both statements compared the retrieval arm
against the \emph{prompt} baseline, which is not a no-context condition because it carries a
developer-relations instruction. Against the arm that sees the review text only, retrieval scores
2.24 against 2.31, $-0.07$ at $p = 0.90$, and all three arms come from the discarded
template-composer round in any case, so none supports a claim in either direction. Stage~4 leaves
a null and nothing else; Stage~3's contribution stands on its own evidence in
Section~\ref{exp1-results} and no longer has a demonstrated downstream payoff.

\subsection{Aligning the IR-Aware Generator at Scale}\label{exp3-results}

  Five policies were trained on distilGPT2 (settings in Appendix~\ref{appendix:settings}).
  Head-to-head on a 100-review test set, BLEU-1~/~ROUGE-L~/~BERTScore~F1 are 0.090~/~0.097~/~0.802
  for SFT-base, 0.068~/~0.074~/~0.792 for KTO, 0.084~/~0.088~/~0.800 for DPO,
  \textbf{0.137~/~0.131~/~0.806} for the constrained proxy, and 0.087~/~0.091~/~0.798 for
  Lagrangian PPO, so the constrained proxy leads three of four metrics with $+53\%$ BLEU-1 over
  SFT-base; the rule-based safety-violation rate is uninformative here because compliance is
  satisfied at initialization.

  \textbf{The dual-objective claim is not demonstrated, and we report that as the result.} At
  distilGPT2 scale the compliance scorer is satisfied at initialization ($C = 0.94 \geq \tau$), so
  $\lambda$ collapses to zero and the update reduces to single-objective optimization. Two further
  experiments show this is a property of the regime rather than of one threshold choice: a re-run
  with the four operational compliance checks and $\tau$ tightened from 0.50 to 0.90 records 0
  binding steps out of 30 and one violation across the run, and a sweep over $\tau \in \{0.0, 0.5,
  0.7, 0.9, 1.0\}$ leaves the pass rate at 100\% for four of the five policies at every threshold,
  only SFT-base dipping to 99\% at $\tau \geq 0.9$. A model of this size cannot plausibly violate
  the rubric, so the CMDP machinery is untested under binding conditions. A second explanation we
  cannot separate from the first weakens the finding further rather than rescuing it. The
  operational scorer is a regular expression over roughly thirty tokens, and we tested its
  sensitivity by hand: nine replies that plainly breach the four rubric dimensions but are phrased
  away from the patterns (``Our engineers have committed to a permanent resolution by Friday'')
  all score a perfect 1.00, and empty output, whitespace, and gibberish score 1.00 too, so a
  policy that collapses to emitting nothing is recorded as perfectly compliant. The quality scorer
  has the mirror defect: eleven rubric keywords in a bare list score 1.000 against 0.672 for a
  genuine reply. Our sweep therefore establishes only that these policies rarely emit about thirty
  specific tokens; whether they violate the four behaviours the rubric names this instrument
  cannot tell us, and we no longer claim it does. Every compliance number here is a property of
  the detector, and the $+53\%$ BLEU-1 figure is a single-objective result, not evidence of a
  quality-versus-compliance trade-off.

  \textbf{The constraint binds once the policy can violate it, as this instrument measures
  violation.} We originally read the obstacle as capability rather than parameter count:
  distilGPT2 never produces the fluent phrasing the rubric penalises, so it cannot violate. Given
  the detector's insensitivity that reading is not the only one, and every violation rate below
  carries that ambiguity. We first justified it by observing that 5.2\% of our own Stage-4
  generations trip the scorer. That figure is from the four withdrawn template-composed arms (21
  violations in 400 replies); scored the same way, the two model-generated arms we actually defend
  violate on 4 of 200 replies, 2.0\%, one of them an information leak rather than an over-promise.
  The zero-shot Qwen2.5-1.5B probe gives 2 violations in 40 replies, 5.0\%, whose 95\% interval of
  roughly 1--17\% cannot be distinguished from 2.0\% either way, and distilGPT2 gives roughly
  zero. The honest version is weaker than the one we first made: a 1.5B instruct model violates
  occasionally where distilGPT2 essentially never does, enough for the constraint to bind, but we
  cannot claim it matches the rate of our own generations. Re-running the Lagrangian loop with
  Qwen2.5-1.5B-Instruct as the policy (Appendix~\ref{appendix:settings}) inverts the behaviour:
  over 90 steps the constraint binds on 13 (14.4\%), and $\lambda$ leaves zero, rises
  monotonically to 0.65 (0.05, then 0.40, then 0.65 by thirds), and never collapses back, while a
  30-step run reaches only $\lambda = 0.125$, so the multiplier scales with the horizon. The claim
  that the CMDP degenerates to single-objective optimization is therefore a property of the
  proof-of-concept policy, not of the formulation.

  \textbf{What is and is not shown.} Quality moves from 0.565 to 0.590 and compliance from 0.997
  to 0.989 while $\lambda$ climbs, but we do not call this a trade-off: the step-level correlation
  between the two is $\rho = -0.12$ ($p = 0.25$), two separate trends over one run with no
  measured coupling. Nor can we show the multiplier growing large enough to pull the policy back
  into compliance; at a 5\% base violation rate and a batch of four, roughly one violating sample
  arrives every seven steps, and 90 LoRA steps do not accumulate enough pressure to reverse the
  drift. So the dual objective engages and the multiplier responds monotonically to the tokens it
  matches, the tension between objectives is not measurable on our data, and closing the loop
  needs a longer horizon or a higher base violation rate than a 1.5B instruct model supplies. What
  this layer contributes is the formulation instantiated over four dev-rel constraints, a trainer
  that runs, the threshold sweep, the capability probe, the binding-constraint re-run, and the
  released testbed.

\subsection{Cluster Quality and Purity}\label{cluster-validation-results}

  Stage~2 was redesigned from a flat UMAP+HDBSCAN baseline (194 clusters, mean size 471.0) to the
   three-layer KG-hierarchical design (605 sub-clusters, mean size 15.9), but the two were not run
   over the same reviews: the flat baseline clusters 91{,}368 reviews from the full relabelled
   corpus, whereas the KG design is built on the 8{,}404-review sample, of which 5{,}689 land in a
   sub-cluster. That sixteen-fold difference is a third reason, on top of the two below, why the
   head-to-head intrinsic comparison is not a comparison.
   We compute intrinsic geometry (Davies-Bouldin~\cite{davies1979cluster},
  Calinski-Harabasz~\cite{calinski1974dendrite}, silhouette~\cite{rousseeuw1987silhouettes},
  following~\cite{schutze2008introduction, steinbach2000comparison, strehl2002cluster}) and
  extrinsic purity~\cite{amigo2009comparison} under the Y/P/N weighted rubric of
  Appendix~\ref{appendix:annotators}.

We reported that KG-hierarchical achieves $5.4\times$ lower Davies-Bouldin and $1.9\times$
  higher Calinski-Harabasz than the flat baseline. We withdraw those two figures. The script that
  produced the flat baseline's intrinsic metrics draws its cluster labels as a random subsample of
  10{,}000 reviews spanning all 194 clusters, but collects the review texts to embed by walking
  the clusters in order, and the first flat cluster alone contains 10{,}507 reviews. The
  embeddings therefore come from one cluster while the labels span 194, so the flat DB of 12.15
  and CH of 0.98 measure random labelling rather than the flat partition, and the comparison was a
  valid measurement (the KG's DB 2.24 and CH 1.85, on a correctly paired set) against a broken
  one. Repairing the alignment would not make the two comparable anyway, since the flat side
  embeds 10{,}000 full reviews and the KG side three representatives per cluster. The count-matched
  A1b ablation, which embeds representatives on both sides, is the only intrinsic comparison we
  report, and it is unaffected. Extrinsic purity is also unaffected: the KG's finer clusters score 0.39 against
  the flat baseline's 0.625 under the same Qwen judge, and the flat baseline's 0.66 lead-author
  purity sits in the substantial-agreement band ($\geq 0.61$) against KG-605's fair band ($\geq
  0.21$) on the Landis-Koch~\cite{landis1977measurement} convention adapted to weighted purity,
  reflecting the granularity-versus-coherence trade-off A1b makes explicit. Per issue type the KG
  loses on four of five.

  \textbf{Where the KG actually wins.} On the prioritization axis the story reverses: of the
  10{,}534 aspect nodes, the top five by PageRank are touched by 55\% of reviews and the top ten
  by 66\%, so a developer triaging in that order reaches 55\% of the complaint volume after five
  aspect buckets. These are the deduplicated figures; counting a review once per aspect it
  mentions inflates them to 74\% and 96\%, and that pair should not be used. We previously wrote
  ``three-quarters'' here, the discarded 74\% reappearing two lines after we told the reader not to
  use it, and also claimed a ``roughly $10^3\times$ concentration ratio'' against random aspect
  ordering. We withdraw that too: we never computed it, and random ordering is the wrong baseline
  anyway, since the operative comparison is ordering by raw frequency, which we did not run. The KG
  additionally enables per-aspect drill-down, distinguishing \texttt{battery}$\to$Samsung-drain
  from \texttt{battery}$\to$fast-charge.

\subsection{Ablations}\label{ablations}

  An ablation removes one component and re-measures. Eight are listed, counting A1b; seven were
  run and A7 is deferred. Three came back against the design, and Table~\ref{tab:evidence} carries
  the verdicts.

  \textbf{A1, remove the knowledge graph and feed flat clusters into Stage~3.} Stage~3 quality
  does not change, because the translator consumes a cluster centroid and its representatives, not
  the graph edges, so the KG is not a Stage~3 dependency.

  \textbf{A1b, compare the KG against flat clustering matched on cluster count.} Intrinsic metrics
  move with cluster count on their own, so comparing the KG's 605 sub-clusters against the flat
  baseline's 194 conflates structure with granularity. Re-running agglomerative clustering forced
  to 605 clusters, on 1{,}815 representative reviews, removes the confound, and the flat version
  wins on all three metrics:
  Davies-Bouldin 1.12 against 2.24 (lower is better), Calinski-Harabasz 4.21 against 1.85 (higher
  is better), and silhouette $+0.148$ against $-0.234$ (a negative value means the average point
  sits closer to a neighbouring cluster than to its own). The KG does not produce geometrically
  better clusters; what it produces is an ordering, and prioritization is the only cluster-layer
  claim we make.

  \textbf{A2, remove hierarchical sub-clustering.} Weighted Y/P/N purity is 0.625 for the flat
  194 and 0.390 for the KG's 605 (Section~\ref{cluster-validation-results}): finer granularity
  costs coherence rather than buying it, the second result against the cluster layer.

  \textbf{A3, remove taxonomy grounding.} The LLM writes a free-form issue description instead of
  filling a routed template. Strict template-fill falls from 0.96 to 0.69 and the per-type
  structural columns fall to exactly zero (Section~\ref{exp1-results}). This is the ablation that
  isolates the paper's central mechanism: the gain is not the language model, it is the routing.

  \textbf{A4, remove the HITL-2 spec review.} Benchmark scores do not move, for a procedural
  reason: HITL-2 is a production-time gate and the 100 evaluation specs were scored before any
  human correction, so the headline numbers already describe pre-HITL output. The ablation says
  nothing about whether HITL-2 helps in deployment, and we do not claim that it does.

  \textbf{A5, remove retrieval and keep the IssueSpec.} BLEU-1 moves by $-0.006$, ROUGE-L by
  $-0.008$, and BERTScore F1 by $-0.004$, all under one percent absolute. Three caveats belong
  with this: A5 was generated by the template composer that Section~\ref{exp2-results} withdraws,
  so it measures retrieval's effect on a generator we no longer report; it was scored on automatic
  metrics only while A6 was scored by humans, so the contrast between them is directional and a
  human-rated no-retrieval condition is the missing cell; and the companion claim that retrieval
  without the IR falls below a no-context baseline is withdrawn (Section~\ref{exp2-results},
  $-0.07$, $p = 0.90$). A5 therefore reads simply as retrieval making no measurable difference
  here, with no accompanying story about it doing harm.

  \textbf{A6, remove the IssueSpec and keep retrieval.} Superseded. Under the confounded design
  it showed quality falling from 4.59 to 2.24; under the corrected same-model design
  (Section~\ref{exp2-results}) removing the IssueSpec changes quality by $-0.04$, not
  distinguishable from zero. A5 and A6 together now say something different from what we first
  concluded: neither retrieval nor the typed intermediate produces a measurable change in the
  reply a user receives, and the Stage~3 to Stage~4 coupling we claimed is not supported.

  \textbf{A7, single-stream RLHF feedback.} Not run. It tests whether quality and compliance
  feedback can share one channel, which is only meaningful once the constraint binds and the two
  signals compete; Section~\ref{exp3-results} reaches the first half of that condition and not the
  second, so A7 stays deferred rather than reported.

%% =====================================================================
\section{Discussion}\label{discussion}

\subsection{What Carries the Result, and What Does Not}\label{load-bearing}

\emph{Load-bearing.} Two components carry the measured gains: type-routed templating, since
removing taxonomy grounding (A3) drops strict template-fill from 0.96 to 0.69 and collapses the
per-field metrics to zero, and the verified-anchor procedure, whose evidence and caveats are in
Section~\ref{exp1-results}. The typed IR is \emph{not} a demonstrated driver of Stage~4 quality:
removing it under a same-model comparison moves quality by $0.04$ ($p = 0.54$,
Section~\ref{exp2-results}), removing retrieval instead (A5) moves every automatic metric by under
one point absolute, and the retrieval-versus-no-context claim we made twice does not survive
either. Nothing about retrieval survives from Stage~4 beyond the null.

\emph{Not load-bearing, and reported as such.} The knowledge graph does not win on intrinsic
cluster geometry: under count matching (A1b) flat agglomerative clustering beats it on all three
metrics, and sub-cluster purity is lower at the KG's finer granularity (A2). Its defensible
contribution is centrality-based prioritization, and even that goes less far than it appears,
since PageRank runs on the bipartite review-aspect graph and nothing shows it needs the
hierarchical sub-clustering layer, a configuration we did not test. Stage~2 as built is therefore
justified in prioritization and unjustified in sub-clustering. Stage~5 sits in between: the
constraint does bind once the policy can violate it and the dual update engages, but no
quality-versus-compliance trade-off is demonstrated at the scale we can train at
(Section~\ref{exp3-results}), so it ships as a testbed with a partial result. \textsc{SpecCov}
falls in the same group for a sharper reason: against human faithfulness judgements it shows no
correlation at all (Section~\ref{speccov-validation}), so we withdraw it.

What survives ablation, then, is the typed IR, the type-routed templating that produces it, and
the verified-anchor procedure that makes it producible at corpus scale; a reader who wants the
smallest system reproducing that needs Stage~1 and Stage~3. Stage~4 is not part of it, and we
state the consequence rather than soften it: on our evidence Stage~4 contributes nothing we can
measure, and a reader building on this work should treat the reply-generation layer as unproven
rather than as a component with a small effect.

\subsection{Implications}
Three claims travel beyond app reviews. \emph{First}, LLM annotation noise on
  software-engineering text is not uniform, but the shape is not the one we first reported. We
  had described the 25\% mislabelling rate as concentrated on minority classes. Recomputing per
  stratum on the verified anchor shows the opposite: 1{,}302 of the 1{,}307 disagreements fall in
  the \texttt{praise} stratum, the largest class, at 25.8\%, against 0.5\% on
  \texttt{performance}. What is true, and narrower, is that the labeller never emits two of its
  seven classes at all across the corpus, 8 \texttt{compatibility} and 184 \texttt{performance}
  labels in 215{,}583 reviews. Systematic absence of a class, rather than elevated error on rare
  classes, is what a small verified anchor is able to repair. \emph{Second}, a caution rather than a claim. We expected a typed structural intermediate to do
  the load-bearing work in spec-aware response generation and reported it as such;
  Section~\ref{exp2-results} withdraws that, and the only downstream movement we can still point
  to is $+0.12$ on a rule-based scorer from agentic iteration over ten reviews, too small to build
  on. What generalises is the negative lesson: a structured intermediate can be demonstrably
  better as an artifact and still leave the generated text indistinguishable, so pipelines of this
  shape should measure the two separately. \emph{Third}, A1b is a self-contained methodological contribution: whenever a clustering design
  changes cluster count and structure at once, intrinsic geometry metrics conflate the two, so the
  minimum reporting standard is a count-matched baseline paired with a domain-utility metric
  (PageRank coverage here).

\subsection{Threats to Validity}
\emph{Construct.} Template-fill measures coverage rather than content correctness, and \textsc{SpecCov} does not fill that gap (Section~\ref{speccov-validation}), being a lexical-grounding proxy that does not track human faithfulness judgement~\cite{grusky2018newsroom, maynez2020faithfulness, kryscinski2020evaluating, ji2023survey}. The paper therefore has no validated content-correctness measure for Stage~3 and the rubric carries that burden alone; NLI- and QA-based measurement~\cite{laban2022summac, honovich2022true, durmus2020feqa} is the next step.
\emph{Internal.} The count-versus-structure confound is exposed and resolved by A1b, which reframes the knowledge graph's contribution as prioritization rather than geometry. A threat we did not remove is that the 100 clusters carrying the Stage~3 and Stage~4 results come from the KG pipeline, the layer A1b and A2 find wanting; since flat clustering scores higher on purity, the headline gains may be measured on a slightly worse partition, which would understate rather than inflate them, and re-running Stages~3 and~4 over flat clusters is the cheapest open experiment in the paper. Two further internal threats are stated with the Stage-1 results: the gold standard is seeded with V5's predicted label and stratified over V5's predicted classes (Section~\ref{experiments}), and every \texttt{compatibility} figure rests on 300 augmented training rows, one figure from which we withdraw outright (Section~\ref{exp1-results}).
\emph{External.} RRGen contains 58 Android apps, named by real package identifier rather than anonymised; V5 transfers zero-shot to the Maalej label space, but Stage~2--4 cross-corpus transfer remains unvalidated, GPT-4o and Gemini-1.5-Pro replication is queued (Section~\ref{future-work}), and every KG-hierarchical number describes a 3.9\% sample of the corpus (Section~\ref{aim1}).
\emph{Conclusion.} The three-rater agreement on the classification gold ($\kappa = 0.968$) sits well above what free-labeling studies report, which invites the question of whether the raters were genuinely independent. Three properties of the design account for the gap: raters verify a supplied label rather than choose one of seven, collapsing most of the decision space; a 20-review calibration round with adjudicated disagreements precedes the main task; and the protocol fixes eight boundary rules. Residual disagreement is rater-specific rather than random, one rater accounting for most \texttt{usability} reassignments, which is the pattern independent annotation produces.
\emph{Reliability, and which coefficient belongs to what.} The three agreement statistics here measure different things. The classification gold reaches Fleiss $\kappa = 0.968$ and $0.979$ across three human raters. The reported Stage-4 set, rated by the same three humans on all 200 rows, has quality weighted $\kappa$ of 0.03, 0.04 and 0.50; that is low and it is the number to quote, because the 0.970 to 0.986 range belongs to the discarded round where one arm was visibly broken. The weakest figure, $\alpha = 0.285$ (corrected from a mis-implemented 0.451, Section~\ref{experiments}), comes from a deliberately weaker LLM-assisted panel~\cite{gilardi2023chatgpt}; no claim rests on it alone.
\emph{Statistical.} Experiment~2 reports a null, so the question is power rather than multiplicity: the design rules out a large effect and not a small one (Section~\ref{exp2-results}). In the other direction, the one cross-family quality difference reaching nominal significance does not survive multiplicity correction and is withdrawn (Section~\ref{exp1-results}).

\subsection{Reproducibility and Extensibility}\label{reproducibility}
Every numerical claim in Section~\ref{results} is backed by a saved data file, and the full path
from raw RRGen to Stage~4 output takes about six hours on one GPU workstation. We audited what a
reader can check without that run by execution rather than assumption. Re-running the released
analysis scripts regenerates roughly sixty stored result files bit-identically, including every
statistical test behind RQ1 and RQ2, but what those scripts consume is not itself regenerable: the
Stage-3 specification sets, the 100-cluster sample, and the Stage-4 RAG sample ship as artifacts
with no standalone producer script, since producing them requires the proprietary-API generation
pass. The analysis layer is therefore reproducible from the middle outward and the data-generation
layer only against a paid API. One artifact does not pin at all: the GitHub-mined comparison set
is re-queried live and returns a different issue mix each run, so the \texttt{human\_github}
condition rests on a snapshot we can release but not regenerate.

The release bundles the V1--V5 checkpoints and cleanlab procedure, the verified anchor and gold,
both Stage-4 rating rounds with the generated replies and the rating-provenance check, the A1b
output, the \textsc{SpecCov} scorer with its 60-spec faithfulness ratings, the CMDP-RLHF testbed,
the 100-cluster benchmark with the LLM-assisted subsample kept separate from the all-human sets, a
pinned Docker image, and a nineteen-segment verifier that recomputes numerical claims from the
saved files, the discarded Stage-4 round retained as a segment labelled withdrawn. We audited the
verifier itself, and its coverage is narrower than we
previously stated: three segments recompute a value from the underlying data and assert it, while
the remainder load a stored summary and print it, which checks that the manuscript agrees with the
released file but not that the file is right. Uncovered numbers include the Stage-3 template-fill
table, the held-out $\kappa$ of 0.616, the Qwen-judge purity figures, the Maalej cross-protocol
row, the per-dimension cells of the cross-family rubric, and most secondary statistics in
Sections~\ref{exp2-results} and~\ref{exp3-results}. We list this rather than repeat the earlier
and much smaller claim that only two tables were uncovered.

Because each stage emits a typed intermediate, components can be substituted independently.
Deterministic components use fixed seeds; the LLM-dependent stages reproduce statistically rather
than bit-identically, and we measured that variance. Regenerating Stage~3 specs three times over
the same 15 stratified clusters with the local Qwen2.5-3B generator at temperature 0.7, changing
only the seed, gives strict template-fill 0.685, 0.700, and 0.634 (mean 0.673, sd 0.029) and
\textsc{SpecCov} 2.73, 2.67, and 3.00 (mean 2.80, sd 0.144), with no parse failures. Run-to-run
noise on template-fill is therefore about three points on a 0--1 scale, and that is the yardstick
for the differences we report: the 0.96 versus 0.69 for taxonomy grounding clears it comfortably,
the withdrawn Stage-4 effect at $+0.04$ sits an order of magnitude below it, and gaps of a few
hundredths such as 4.09 versus 3.96 in the cross-family rubric do not clear it. We previously also
cited the 0.96-versus-0.48 cross-family span here; that comparison is not matched on sample or
issue type (Table~\ref{tab:multi-llm}), so we no longer use it as a magnitude. Two limits remain:
neither the headline Claude generations nor Stage~4 response generation was repeated, because that
needs API budget rather than new code, and the released script runs both sweeps unchanged.

\subsection{Future Work}\label{future-work}
Five extensions start from released resources.
\textbf{(0)}~Replace \textsc{SpecCov} with an entailment-based faithfulness
measure~\cite{laban2022summac, honovich2022true, durmus2020feqa}, validated against the 60-spec
human ratings we release, and widen the cross-family rubric comparison beyond the 14 matched
bug-report clusters with a judge outside all four families.
\textbf{(1)}~Ask the right raters before scaling: a triage reply is consumed by a developer, so
the next Stage-4 round should recruit raters with software-engineering experience and no stake in
the system before growing the sample to 200--300 clusters.
\textbf{(2)}~Run the CMDP loop where the trade-off could show, with a longer horizon and a
higher base violation rate on an 8B-class model.
\textbf{(3)}~Frontier-family replication (GPT-4o, Gemini-1.5-Pro) on Stages~3--4, and
\textbf{(4)}~cross-corpus transfer of Stages~2--4, V5 already transferring zero-shot to
Maalej~\cite{maalej2015bug}.
\textbf{(5)}~An agentic planner over Stage~4 dispatching a different sub-agent profile per issue
type, using SWE-Agent~\cite{yang2024swe}, RepairAgent~\cite{bouzenia2025repairagent}, or
HyperAgent~\cite{phan2024hyperagent} for defect repair; a spec-level dispatch stub already runs
across all five issue types, and end-to-end patch evaluation is the next step.

%% =====================================================================
\section{Conclusion}\label{conclusion}

  We proposed \textbf{IssueSpec}, a knowledge-graph-grounded typed intermediate representation
  aimed at the five gaps of Section~\ref{introduction}, closing the first three and reporting
  negative results against the last two. A verified-anchor procedure makes the IR producible at
  corpus scale from noisy LLM labels, lifting Cohen's $\kappa$ from 0.163 to 0.592, and to 0.616
  on a held-out subset whose value is replay-dependent and indistinguishable from the full-set
  figure; classifier endorsement of the corrections is 88.66\%, a consistency check rather than
  independent validation. Cross-family replication on four LLMs preserves the rank order
  descriptively, with no pairwise difference significant, so it is weak evidence against
  single-model bias rather than a ruling-out. We set out to show that the typed IR is the
  load-bearing component of spec-aware response generation, and we cannot: holding the generator
  fixed moves human-rated quality by $+0.04$ Likert points ($p = 0.54$, three raters, 100 paired
  reviews) with the raters disagreeing in direction, and the $+2.35$ we previously reported came
  from a comparison whose arms were written by different template composers. Nor does the negative
  half of our original hypothesis survive: we had reported that retrieval alone falls below a
  no-context baseline, and against the arm that actually has no context it scores 2.24 against
  2.31 at $p = 0.90$, inside the discarded round in any case. Stage~4 yields one usable result,
  the null. Two findings travel beyond app reviews: an LLM
  labeller on software-engineering text can omit whole classes rather than merely err on them, and
  a structured intermediate can be demonstrably better as an artifact while leaving the text
  generated from it indistinguishable, a distinction pipeline evaluations of this shape should
  measure separately. The released artifacts are listed in Section~\ref{reproducibility}.

%% =====================================================================
\section*{CRediT authorship contribution statement}
\textbf{Fabiha Jalal:} Conceptualization, Methodology, Software, Validation, Formal analysis,
Investigation, Data curation, Writing -- original draft, Writing -- review and editing,
Visualization. \textbf{Sadia Tasnim Dhruba:} Investigation, Data curation, Validation, Writing
-- review and editing. \textbf{Tazrian Zaman Tushti:} Investigation, Data curation, Validation,
Writing -- review and editing. \textbf{Ajwad Abrar:} Methodology, Supervision, Writing --
review and editing. \textbf{Md Kamrul Hasan:} Supervision, Resources, Writing -- review and
editing. \textbf{Hasan Mahmud:} Supervision, Resources, Project administration, Writing --
review and editing.

\section*{Declaration of competing interest}
The authors declare that they have no known competing financial interests or personal
relationships that could have appeared to influence the work reported in this paper.

\section*{Data availability}
All code, configuration files, and analysis scripts are available in the project repository.
A verification bundle that recomputes or reproduces the numerical claims in
Section~\ref{results}, with the coverage limits stated in Section~\ref{reproducibility}, is archived on Zenodo
(DOI 10.5281/zenodo.22048738, v3), and the Stage-1 V5 classifier is released on Hugging Face.
That bundle supersedes two earlier releases (DOI 10.5281/zenodo.21982774 and DOI 10.5281/zenodo.20320410), the first of which redistributed corpus-scale review text we had no right to redistribute and the second of which also predates the
Stage-4 correction described in Section~\ref{exp2-results}; both versions remain resolvable, and
the earlier one should not be used.
The underlying review corpus derives from the RRGen dataset of Gao et al. That release is
access-gated behind a request form, states academic use only, and carries no licence granting
redistribution, so we do not redistribute it. We also do not redistribute the Maalej and Nabil or
Guzman corpora. An earlier version of our repository did redistribute all three; we removed them,
and \texttt{DATA\_NOTICE.md} records what was removed and how to obtain each corpus from its
original authors. Derived artifacts that embed verbatim review text carry the same restriction as
the source and are not offered under an open licence; original code is released under MIT.

\section*{Ethics, consent, and data protection}
\emph{Human participants.} Three people annotated for this study: the lead author and two
collaborators. The annotation protocol, released with the artifact, states that participation was
voluntary, that annotators could withdraw at any time, that no personal information was collected
beyond a name and email for attribution, and that compensation took the form of authorship credit.
Both additional annotators are consequently co-authors of this paper, which we state explicitly
because the paper's discussion of rater independence should be read with it in mind. The work was
not submitted to an institutional review board. Under our institution's practice a study of this
kind, involving colleagues rating text with no collection of sensitive personal data, does not
require review; we record that as our reading rather than as a formal exemption, since no
exemption number was issued.

\emph{Review authors.} The corpora we work with are user-generated public reviews whose authors
did not consent to research use and cannot practically be contacted. We took three positions on
this. We do not redistribute any of the three source corpora, for the licensing reasons in the
data-availability statement and because the Maalej and Nabil files contain unredacted email
addresses, telephone numbers, and personal names. Derived files we do release contain review text
as RRGen supplies it, lowercased and lemmatised with digits masked; that masking is imperfect, and
we found a small number of surviving transaction identifiers in the source, which anyone
re-releasing derived text should re-scan for. We report no result at the level of an individual
review author, and no reviewer identifier appears in any released file.

\emph{Machine-generated replies.} The replies this system produces read as human-written, so any
deployment should label them as machine-assisted, which we regard as a condition of use rather
than a recommendation.

\section*{Role of the funding source}
This research did not receive any specific grant from funding agencies in the public,
commercial, or not-for-profit sectors.

\section*{Acknowledgements}
We thank the human annotators who produced the verified anchor and the lead-author reference
set. The authors verified all technical content, experimental results, and claims in this paper.

%% =====================================================================
\bibliographystyle{elsarticle-num}
\section*{Declaration of Generative AI and AI-assisted technologies in the writing process}
We disclose all uses of generative AI in this work in line with Elsevier policy.
\emph{Writing and code assistance.} An earlier version of this disclosure said that Claude
assisted only with grammatical refinement, caption polishing, and LaTeX formatting, and that all
pipeline code, experimental scripts, and analysis were written by the human authors. That
understated the use, and the repository we release contradicts it: at the time of writing 49 of
58 commits (84\%) carry a \texttt{Co-Authored-By: Claude} trailer, including commits that add pipeline stages, the
constrained-PPO trainer, the analysis scripts, and the correction passes described in
Section~\ref{exp2-results}. The accurate statement is that Anthropic's Claude was used
throughout implementation and analysis as a coding assistant, under author direction and review,
and also for manuscript editing. We prefer to correct this here rather than leave a disclosure
that the artifact refutes. Every structural, framing, experimental-design, and claim decision was
made by the authors, and every reported number was produced by author-written code run under
human supervision and verified against saved data files. \emph{Models inside the system.} Several
generative models are pipeline components and are reported as such in the body
(Table~\ref{tab:multi-llm}): Claude Opus 4.7 as the headline Stage~3 generator, Llama-3.3-70B and
Qwen2.5-\{3B, 1.5B\} as cross-family comparators, and Qwen2.5-3B as the LLM-as-judge for the
rubric and the purity audit, both calibrated against the lead author. The V2 corpus labels are
LLM-generated and explicitly treated as noisy, which motivates the verified-anchor procedure of
Stage~1. \emph{Human verification.} No reported numerical result, expert annotation, or ground-truth label was
automatically generated by AI without human verification. The 5{,}230-review verified anchor, the
490-review classification gold, the 99-review three-rater ratings, and the 20 lead-author
IssueSpec references were all produced and reviewed by human annotators
(Section~\ref{experiments}).

%% =====================================================================

\bibliography{references}

\appendix

  \section{Implementation Settings}\label{appendix:settings}

  Section~\ref{system-design} states what each stage does and why; this appendix collects the
  settings. Stage~1 fine-tunes a RoBERTa multi-label head; confident-learning correction against
  the verified anchor is gated by anchor confidence $\geq 0.70$ and LLM probability $\leq 0.20$,
  with the $3\times3$ sensitivity grid over $(0.65, 0.70, 0.75) \times (0.15, 0.20, 0.25)$ reported
  in Section~\ref{exp1-results}, and the \texttt{compatibility} stratum is seeded as described in
  Section~\ref{aim1}. Stage~2 embeds with Sentence-BERT and clusters with HDBSCAN
  aspect-first over the 8{,}404-review sample with no dimensionality reduction, against a flat
  UMAP+HDBSCAN baseline over 91{,}368 reviews from the full corpus, so the two describe different
  review sets; PageRank runs
  on the bipartite review-aspect graph. Stage~3 routes to one of five templates and re-prompts
  once with dimension-level feedback on any rubric failure.
  Stage~4 retrieves over a 15{,}100-document ChromaDB index built from five sources, with the
  IssueSpec injected per query; the agentic-RAG feasibility study uses 10 stratified reviews, two
  per issue type, with Qwen2.5-3B as critic and reviser and up to two iterations. Stage~5 trains
  five policies on distilGPT2 (82M parameters) from 400 SFT samples, 296 KTO pairs, and 100 DPO
  pairs; the Lagrangian loop runs 30 steps at batch size 4 with KL coefficient 0.05, and the
  compliance threshold sweep covers $\tau \in \{0.0, 0.5, 0.7, 0.9, 1.0\}$. The binding-constraint
  re-run uses Qwen2.5-1.5B-Instruct with LoRA adapters over the attention projections, the frozen
  base as KL reference, $\tau = 1.0$, and 90 steps, with scorers, dual-ascent rule, and batch size
  unchanged. \textsc{SpecCov} counts substantive-token overlap with the source cluster, tokens of
  length $\geq 5$ after stop-word filtering, with quoted-string and coverage bonuses.

\section{Annotator Guidelines}\label{appendix:annotators}

  Section~\ref{experiments} describes who annotated what; this appendix records the instruments.
  Anchor and gold annotation assign one of seven classes (\texttt{bug\_report},
  \texttt{feature\_request}, \texttt{performance}, \texttt{usability}, \texttt{compatibility},
  \texttt{praise}, \texttt{other}), multi-label permitted, with items below a ``fairly confident''
  threshold skipped. The IssueSpec rubric scores each spec 1--5 and independently per dimension on
  \emph{completeness} (are all required template fields filled?), \emph{accuracy} (do field values
  reflect the source cluster?), \emph{actionability} (would a developer know what to do?),
  \emph{specificity} (does the spec name concrete entities?), and \emph{clarity}. The Stage~4
  instrument is \emph{quality} (1--5), \emph{specificity} (1--5), and \emph{helpful} (Y/N). The
  Y/P/N cluster-purity rubric reads 5 representative reviews per cluster and assigns \emph{Y} (all
  5 share the sub-theme), \emph{P} (3 or 4 share), or \emph{N} (fewer than 3), with weighted
  purity $\mathrm{purity}_w = (|Y| + 0.5|P|)/(|Y|+|P|+|N|)$. All comparative evaluations were
  blinded throughout the rating window, condition labels replaced by random codes and unsealed
  only at aggregation.

\end{document}
